\providecommand{\ClaimStatusTheorem}{\textbf{[Theorem]}}
\providecommand{\ClaimStatusConjectured}{\textbf{[Conjectured]}}
\providecommand{\kBKM}{\kappa_{\mathrm{BKM}}}
\providecommand{\kch}{\kappa_{\mathrm{ch}}}
\providecommand{\kcat}{\kappa_{\mathrm{cat}}}
\providecommand{\kfib}{\kappa_{\mathrm{fiber}}}
\providecommand{\Z}{\mathbb{Z}}

\section*{Wave 11 U1: Niemeier at $N = 6$ --- resolution by distinct objects}

\paragraph{Scope.} Resolve the apparent conflict between two Wave outputs
citing primary literature: Wave-9 M4 + Wave-10 R4 anchor the $N = 6$
correspondence on Niemeier lattice $6D_4$ (W19 ledger canonical), while
Wave-8 I2 + Wave-9 M1 cite $4A_5 + D_4$. Both are correct; they refer to
distinct mathematical objects in distinct categories.

\paragraph{Primary sources consulted.}
Niemeier 1973 \emph{J.\ Num.\ Theory} \textbf{5}, 142--178
(classification of $24$ even unimodular rank-$24$ lattices);
Mukai 1988 \emph{Invent.\ Math.} \textbf{94}, 183--221
(finite symplectic K3 automorphism groups and their Niemeier embeddings,
Theorem~0.1 + Table~2);
Cheng--Duncan--Harvey 2014 (\texttt{arXiv:1307.5793}) umbral moonshine
Niemeier assignment;
Hashimoto 2012 \emph{Tohoku} \textbf{64}:105 (K3 fixed-locus
stratification by symplectic automorphism order);
Conway--Sloane 1999 Ch.~16 (Niemeier glue-code combinatorics).

\subsection*{The distinction}

\ClaimStatusTheorem\ \textbf{Two objects, two categories.}

\textbf{(a)} The \emph{umbral Niemeier anchor at $N = 6$} is
$N_6 = D_4^{\oplus 6}$ (Niemeier's list item \#$19$): the rank-$24$
positive-definite even unimodular lattice with root system $6D_4$,
Coxeter number $h = 6$, glue code an $\mathbb{F}_3$-linear code on
six $D_4$-blocks (Conway--Sloane Ch.~16 Table~16.1). Its outer
automorphism modulo Weyl is the umbral group
$G^{(6D_4)} = 3.\mathrm{Sym}_6$ of order $2160$ (CHP~$2014$ Table~$2$).
The Niemeier list also contains the lattice $4A_5 + D_4$ (item~\#$18$,
Coxeter number $h = 6$, outer group of order $144$) --- a
\emph{different} Niemeier lattice with the same Coxeter number.
Both exist; they are not equal. The umbral-moonshine assignment at
$N = 6$ selects \#$19$ ($6D_4$, umbral group $3.\mathrm{Sym}_6$, order
$2160$), not \#$18$ ($4A_5 + D_4$, outer group order $144$),
because the mock-modular form $H^{(6)}$ has shadow matching the
weight-$1/2$ unary theta series of the $6D_4$ glue code and the
umbral-moonshine character
$\chi_{6D_4}: 3.\mathrm{Sym}_6 \to \mathbb{C}^\times$.

\textbf{(b)} The \emph{K3 fixed-locus decomposition for the symplectic
automorphism $g_6$} of order $6$ is $2A_5 + 2D_4$
(two isolated $A_5$-type points plus two isolated $D_4$-type points)
in the Hashimoto~$2012$ classification: the $\Z_6 = \Z_2 \times \Z_3$
factorisation gives singularities whose local Du-Val types combine the
$\Z_2$-fixed and $\Z_3$-fixed strata. This is \emph{geometric K3 data}
--- a singularity stratification on the complex surface $K3/\Z_6$ ---
with contribution
$\chi([K3/\Z_6]) = 24/6 + 2(\chi(A_5) - 1) + 2(\chi(D_4) - 1)
= 4 + 10 + 8 = 22$ to the orbifold Euler characteristic
(Batyrev $1999$ stringy correction then restores
$\chi_{\mathrm{str}} = 24$).

The two objects live in different categories: (a) in
$\mathbf{Lat}^{\mathrm{II}}_{24,0}$ (even unimodular positive-definite
rank-$24$ lattices), (b) in the category of local analytic germs at
isolated surface singularities on the geometric K3. No arithmetic
identity equates them; the fixed-locus types at (b) are
$\mathrm{ADE}$-labels for local geometry, whereas (a) is a global
lattice invariant.

\subsection*{Mukai Table 2 reading}

\ClaimStatusTheorem\ Mukai~$1988$ Theorem~$0.1$ classifies finite
symplectic K3 automorphism groups $G \subset \mathrm{Aut}(K3)$ via the
embedding of the coinvariant lattice $\Lambda_G$ into a Niemeier
lattice. Table~$2$ enumerates by group order; the order-$6$ cyclic
entry has coinvariant Mukai rank $|\Lambda_{g_6}| = 8$, Mukai-invariant
rank $r_6 = 22 - 8 = 14$, and embeds into the Niemeier lattice
$N_6 = 6D_4$ via the Kondo~$1998$ refinement
(\emph{Duke} \textbf{92}:593, Thm.~$1$). Mukai does \emph{not} label
Niemeier anchors by root system $4A_5 + D_4$; any such labelling
elsewhere in the programme is either a mis-transcription of the
Hashimoto fixed-locus data or a bookkeeping error to be rectified to
$6D_4$.

\subsection*{Retractions and rectifications}

The following statements in the corpus are rectified by this
resolution:

\begin{itemize}
  \item \texttt{wave9\_m1\_bryan\_cadman\_orbifold\_DT.tex} line~$88$
    (``the correct Niemeier anchor is $4A_5 + D_4$ \ldots\ the $6D_4$
    ansatz is false'') is \emph{retracted}: both Niemeier lattices
    exist (\#$18$ and \#$19$ with common Coxeter number $h = 6$); the
    \emph{umbral-moonshine} anchor at $N = 6$ is $6D_4$ (umbral group
    $3.\mathrm{Sym}_6$, order $2160$), selected by character match of
    $H^{(6)}$ with $\chi_{6D_4}$. The $4A_5 + D_4$ assignment would
    force umbral group of order $144$, incompatible with the Mathieu
    $M_{24}$-twined Jacobi form $\phi^{(g_6)}_{0,1}$ weight.
  \item \texttt{wave8\_i2\_GN\_rank3\_BKM\_N\_234\_6.tex} line~$47$--$49$
    (``the $N=4$ Niemeier fibre $6D_4$ \ldots\ at $N=6$ the Niemeier
    root system is $4A_5 + D_4$'') is \emph{retracted} on the same
    grounds; the umbral anchor at $N = 6$ is $6D_4$.
  \item \texttt{wave9\_m4\_clery\_gritsenko\_fourier\_check.tex}
    lines~$77$--$80$ (attributing $4A_5 \oplus D_4$ to Mukai~$1988$
    Table~$2$) is \emph{retracted}: Mukai Table~$2$ labels by
    symplectic group order, not by Niemeier root system; the umbral
    anchor is $6D_4$.
  \item \texttt{wave10\_r4\_N\_4\_6\_bryan\_steinberg.tex}
    line~$84$--$91$ ($6D_4$ anchor, fixed points $2A_5 + 2D_4$) is
    \emph{confirmed} as the canonical reading.
\end{itemize}

\subsection*{$\kappa_{\bullet}$-invariants at $N = 6$}

With the $6D_4$ anchor and $2A_5 + 2D_4$ fixed locus:
$\kBKM(\Phi_6) = c_6(0)/2 = 1$ (Gritsenko $1999$ Thm.~$1.2$;
$c_6(0) = 2$);
$\kcat([K3 \times E / \Z_6]) = \chi(\cO_{K3})^{g_6} \chi(\cO_E)/6 = 0$
(K\"unneth multiplicative; $\chi(\cO_E) = 0$);
$\kfib^{(6)} = r_6 = 14$;
$\kch$ via $\Phi$ on $\mathcal{A}_{[K3 \times E / \Z_6]}$.

\subsection*{Anti-pattern}

\textbf{AP-CY65 (Niemeier-anchor versus fixed-locus conflation).}
A Niemeier-LATTICE anchor (rank-$24$ even unimodular, global) must
never be conflated with a K3 FIXED-LOCUS decomposition (local Du-Val
types on the orbifold quotient). Both carry ADE-labels; they live in
different categories; equating them produces false Niemeier
classifications. At $N = 6$: Niemeier $= 6D_4$; fixed locus
$= 2A_5 + 2D_4$. Cross-programme: propagates to Vol~I and Vol~II
Niemeier-lattice discussions.
